
\documentclass[aps,preprint]{revtex4}
%%%%%%%%%%%%%%%%%%%%%%%%%%%%%%%%%%%%%%%%%%%%%%%%%%%%%%%%%%%%%%%%%%%%%%%%%%%%%%%%%%%%%%%%%%%%%%%%%%%%%%%%%%%%%%%%%%%%%%%%%%%%
\usepackage{amsfonts}
\usepackage{amsmath}
\usepackage{amssymb}
\usepackage{graphicx}

\setcounter{MaxMatrixCols}{10}
%TCIDATA{OutputFilter=LATEX.DLL}
%TCIDATA{Version=4.00.0.2321}
%TCIDATA{Created=Saturday, May 25, 2002 11:48:07}
%TCIDATA{LastRevised=Thursday, May 30, 2002 13:41:32}
%TCIDATA{<META NAME="GraphicsSave" CONTENT="32">}
%TCIDATA{<META NAME="DocumentShell" CONTENT="Articles\SW\REVTeX 4">}
%TCIDATA{Language=American English}
%TCIDATA{CSTFile=revtex4.cst}

\newtheorem{theorem}{Theorem}
\newtheorem{acknowledgement}[theorem]{Acknowledgement}
\newtheorem{algorithm}[theorem]{Algorithm}
\newtheorem{axiom}[theorem]{Axiom}
\newtheorem{claim}[theorem]{Claim}
\newtheorem{conclusion}[theorem]{Conclusion}
\newtheorem{condition}[theorem]{Condition}
\newtheorem{conjecture}[theorem]{Conjecture}
\newtheorem{corollary}[theorem]{Corollary}
\newtheorem{criterion}[theorem]{Criterion}
\newtheorem{definition}[theorem]{Definition}
\newtheorem{example}[theorem]{Example}
\newtheorem{exercise}[theorem]{Exercise}
\newtheorem{lemma}[theorem]{Lemma}
\newtheorem{notation}[theorem]{Notation}
\newtheorem{problem}[theorem]{Problem}
\newtheorem{proposition}[theorem]{Proposition}
\newtheorem{remark}[theorem]{Remark}
\newtheorem{solution}[theorem]{Solution}
\newtheorem{summary}[theorem]{Summary}
\newenvironment{proof}[1][Proof]{\noindent\textbf{#1.} }{\ \rule{0.5em}{0.5em}}
\input{tcilatex}

\begin{document}

\preprint{}
\title[H paths]{Hilbert Space Paths}
\author{Brian Weiner}

\begin{abstract}
Shell document for REV\TeX
\end{abstract}

\date{05/29/02}
\startpage{1}
\maketitle
\tableofcontents

\section{Paths in Hilbert Space}

Consider the set of paths 
\begin{align}
F& :\bar{\Upsilon}\rightarrow \mathcal{H} \\
F\left( t\right) & =\left\vert \Psi \left( t\right) \right\rangle ;\quad
t\in \tau ,\,\Upsilon =\left( t_{i},t_{f}\right) ;\,\left\vert \Psi \left(
t\right) \right\rangle \in \mathcal{H}.
\end{align}%
Let these paths be strongly continuous on $\bar{\Upsilon}$ i.e. 
\begin{equation}
\lim_{t_{k}\rightarrow t}\left\vert \Psi \left( t_{k}\right) \right\rangle
=\left\vert \Psi \left( t\right) \right\rangle \,\forall \text{ sequence }%
\left\{ t_{k}\right\} \text{ converging to }t
\end{equation}%
i.e.%
\begin{equation}
\lim_{t_{k}\rightarrow t}\left\Vert \left\vert \Psi \left( t_{k}\right)
\right\rangle -\left\vert \Psi \left( t\right) \right\rangle \right\Vert
=0\,\forall \text{ sequence }\left\{ t_{k}\right\} \text{ converging to }t
\end{equation}%
If $F_{1}$ and $F_{2}$ are two such functions then $F_{3}=\alpha F_{1}+\beta
F_{2}$, $\alpha ,\beta \in \mathbb{C}$ is also, as 
\begin{align}
\lim_{t_{k}\rightarrow t}\left\vert \alpha \Psi _{1}\left( t_{k}\right)
\right\rangle & =\left\vert \alpha \Psi _{1}\left( t\right) \right\rangle
;\quad \lim_{t_{k}\rightarrow t}\left\vert \beta \Psi _{2}\left(
t_{k}\right) \right\rangle =\left\vert \beta \Psi _{2}\left( t\right)
\right\rangle   \notag \\
& \Rightarrow \lim_{t_{k}\rightarrow t}\left\{ \left\vert \alpha \Psi
_{1}\left( t_{k}\right) \right\rangle +\left\vert \beta \Psi _{2}\left(
t_{k}\right) \right\rangle \right\} =\left\vert \alpha \Psi _{1}\left(
t\right) \right\rangle +\left\vert \beta \Psi _{2}\left( t\right)
\right\rangle 
\end{align}%
These paths form a linear space $\mathcal{P}\left( \tau ,\mathcal{H},%
\mathcal{H}\right) $ under the sup norm 
\begin{equation}
\left\vert F\right\vert =\sup_{t\in \bar{\Upsilon}}\left\Vert F\left(
t\right) \right\Vert 
\end{equation}%
Subspaces of paths $\mathcal{P}\left( \tau ,\left\vert \Psi
_{i}\right\rangle ,\left\vert \Psi _{f}\right\rangle \right) $ are formed by
paths that have fixed end point values i.e. 
\begin{equation}
\mathcal{P}\left( \tau ,\left\vert \Psi _{i}\right\rangle ,\left\vert \Psi
_{f}\right\rangle \right) =\left\{ \left\vert \Psi \left( t\right)
\right\rangle \in \mathcal{P}\left( \tau ,\mathcal{H},\mathcal{H}\right)
|\left\vert \Psi \left( t_{i}\right) \right\rangle =\left\vert \Psi
_{i}\right\rangle ,\left\vert \Psi \left( t_{f}\right) \right\rangle
=\left\vert \Psi _{f}\right\rangle \right\} 
\end{equation}%
We can generalize this notation to $\mathcal{P}\left( \tau ,\mathcal{M}_{i},%
\mathcal{M}_{f}\right) $ where the paths now have values in the subsets $%
\mathcal{M}_{i}$ and $\mathcal{M}_{f}$ of $\mathcal{H}$, which could be
manifolds. If the manifolds are linear $\mathcal{P}\left( \tau ,\mathcal{M}%
_{i},\mathcal{M}_{f}\right) $ is still a linear space, otherwise this
property is not maintained and $\mathcal{P}\left( \tau ,\mathcal{M}_{i},%
\mathcal{M}_{f}\right) $ becomes a non linear manifold.

By considering the paths 
\begin{equation}
\left\vert e_{j}\left( t\right) \right\rangle =c\left( t\right) \left\vert
e_{j}\right\rangle
\end{equation}
for $j=1,\ldots,\infty$ where $\left\{ \left\vert e_{j}\right\rangle |1\leq
j\leq\infty\right\} $ is a basis for $\mathcal{H}$ one can see that the
values of the paths belonging to $\mathcal{P}\left( \tau,\left\vert \Psi
_{i}\right\rangle ,\left\vert \Psi_{f}\right\rangle \right) $ for any $%
t\in\Upsilon$ is dense in $\mathcal{H}$. In fact as $F$ ranges through all
elements of $\mathcal{P}\left( \tau,\left\vert \Psi_{i}\right\rangle
,\left\vert \Psi_{f}\right\rangle \right) ,$ for a fixed $t\in\Upsilon,$ all
elements of $\mathcal{H}$ are obtained i.e. for a fixed $t$ every element of 
$\mathcal{H}$ lies on some path $\left\vert \Psi\left( t\right)
\right\rangle $.

The space of all paths $\mathcal{P}\left( \infty ,\mathcal{H},\mathcal{H}%
\right) $ contains the subsets $\mathcal{P}\left( \tau ,\mathcal{M}_{i},%
\mathcal{M}_{f}\right) $ where $\mathcal{M}_{i}$ and $\mathcal{M}_{f}$ can
be manifolds or single points. The norm on $\mathcal{P}\left( \infty ,%
\mathcal{H},\mathcal{H}\right) $, (which is a linear space), is given by 
\begin{equation}
\left\vert F\right\vert =\sup_{\Upsilon }\sup_{t\in \bar{\Upsilon}%
}\left\Vert F\left( t\right) \right\Vert 
\end{equation}%
The space $\mathcal{P}\left( \infty ,\mathcal{H},\mathcal{H}\right) $ can be
classified into the linear subsapces $\left\{ \mathcal{P}\left( \infty
,\left\vert a\right\rangle ,\left\vert b\right\rangle \right) |\left\vert
a\right\rangle \in \mathcal{H},\left\vert b\right\rangle \in \mathcal{H}%
\right\} $ and more finely classified into linear subspaces of paths of
different durations 
\begin{equation}
\left\{ \mathcal{P}\left( \tau ,\left\vert a\right\rangle ,\left\vert
b\right\rangle \right) |\left\vert a\right\rangle \in \mathcal{H},\left\vert
b\right\rangle \in \mathcal{H},\left\{ 0\right\} \subseteq \cdots \left[
0,\tau \right] \cdots \subseteq \mathbb{R}\right\} 
\end{equation}

Paths can be varied in the following ways:

\begin{enumerate}
\item within a fixed $\mathcal{P}\left( \tau ,\left\vert \Psi
_{i}\right\rangle ,\left\vert \Psi _{f}\right\rangle \right) $, then the
incremental functions $\left\vert \Delta \left( t\right) \right\rangle $
must belong to the subspace $\mathcal{P}\left( \tau ,0,0\right) $, and thus
periodic in the time duration $\tau $,

\item within the spaces with variable end points, or having one fixed end
point, so that the increment functions $\in $ $\mathcal{P}\left( \tau ,%
\mathcal{H},\mathcal{H}\right) ,$ $\mathcal{P}\left( \tau ,0,\mathcal{H}%
\right) $, $\mathcal{P}\left( \tau ,\mathcal{H},0\right) $.

\qquad Varying the end points is equivalent to varying the paths e.g.%
\begin{align}
\left\vert \Psi \left( t\right) \right\rangle & \in \mathcal{P}\left( \tau
,\left\vert a\right\rangle ,\left\vert b\right\rangle \right)   \notag \\
\left\vert \Phi \left( t\right) \right\rangle & \in \mathcal{P}\left( \tau
,\left\vert a\right\rangle ,\left\vert b\right\rangle +\left\vert \delta
b\right\rangle \right) 
\end{align}%
then define 
\begin{equation}
\left\vert \Delta \left( t\right) \right\rangle =\left\vert \Phi \left(
t\right) \right\rangle -\left\vert \Psi \left( t\right) \right\rangle 
\end{equation}%
which gives that 
\begin{align}
\left\vert \Delta \left( 0\right) \right\rangle & =0  \notag \\
\left\vert \Delta \left( \tau \right) \right\rangle & =\left\vert \delta
b\right\rangle 
\end{align}%
and that $\Delta \in \mathcal{P}\left( \tau ,0,\left\vert \delta
b\right\rangle \right) $. $\Delta $ is a continuous function as it is the
difference of two continuous functions. We can thus obtain all of $\mathcal{P%
}\left( \tau ,\left\vert a\right\rangle ,\left\vert b\right\rangle
+\left\vert \delta b\right\rangle \right) $ by using the single increment
function $\left\vert \Delta \left( \tau \right) \right\rangle $ then using
increment functions belonging to $\mathcal{P}\left( \tau ,0,0\right) $

\item within spaces where the end point are constrained to lie in a proper
sub manifold of $\mathcal{H}$, so that the increment functions $\in \mathcal{%
P}\left( \tau ,\mathcal{M}_{i},\mathcal{M}_{f}\right) $. Clearly the
construction in type 2 only works if $\mathcal{M}_{f}$ is a linear manifold.
If it is not then $\left\vert \delta b\right\rangle $ must belong to the
tangent space of $\mathcal{M}_{f}$.

\item within spaces where the end point are constrained to be fixed but the
time duration varies. Let $\left\vert \Psi \left( t\right) \right\rangle \in 
\mathcal{P}\left( \tau ,\left\vert a\right\rangle ,\left\vert b\right\rangle
\right) $ and $\left\vert \Phi \left( t\right) \right\rangle \in \mathcal{P}%
\left( \tau +\delta \tau ,\left\vert a\right\rangle ,\left\vert
b\right\rangle \right) $ then we can set 
\begin{align}
\left\vert \Phi \left( t\right) \right\rangle & =\left\vert \Psi \left(
t\right) \right\rangle -\left\vert \dot{\Psi}\left( t\right) \right\rangle
\delta \tau \quad \forall t\in \left[ 0,\tau \right]   \notag \\
\left\vert \Phi \left( t+\delta \tau \right) \right\rangle & =\left\vert
\Psi \left( t\right) \right\rangle 
\end{align}%
(the derivative at the boundary is taken to be a one sided derivative). The
incremental paths are then defined as 
\begin{equation}
\left\vert \Delta \left( t\right) \right\rangle =\left\vert \Phi \left(
t\right) \right\rangle -\left\vert \Psi \left( t\right) \right\rangle
;\qquad t\in \left[ 0,\tau \right] 
\end{equation}%
leading to 
\begin{align}
\left\vert \Delta \left( t\right) \right\rangle & =-\left\vert \dot{\Psi}%
\left( t\right) \right\rangle \delta t;\qquad t\in \left[ 0,\tau \right]  
\notag \\
\left\vert \Delta \left( \tau +\delta \tau \right) \right\rangle &
=\left\vert \Psi \left( \tau \right) \right\rangle 
\end{align}%
This construction assumes we use a $\left\vert \Psi \left( t\right)
\right\rangle $ that is continuously differentiable.

\qquad If we start with $\left\vert x\left( t\right) \right\rangle \in 
\mathcal{P}\left( \tau ,\left\vert a\right\rangle ,\left\vert b\right\rangle
\right) $ then we can define the increment as above and obtain $\left\vert
x_{1}\left( t\right) \right\rangle \in \mathcal{P}\left( \tau +\delta \tau
,\left\vert a\right\rangle ,\left\vert b\right\rangle \right) $, then we can
define 
\begin{align}
& \left\vert \delta x\left( 0\right) \right\rangle =0  \notag \\
& \left\vert \delta x\left( t\right) \right\rangle =-\left\vert \dot{x}%
_{1}\left( t\right) \right\rangle \delta \tau =-\left\vert x\left( t\right) -%
\dot{x}\left( t\right) \delta \tau \right\rangle \delta \tau =-\left\vert
x\left( t\right) \delta \tau -\dot{x}\left( t\right) \left( \delta \tau
\right) ^{2}\right\rangle ;  \notag \\
& \left. {}\right. \qquad \qquad \qquad \qquad \qquad \qquad \qquad \qquad
\qquad \qquad \qquad t\in \left( 0,\tau +\delta \tau \right]  \\
& \left\vert \delta x\left( \tau +2\delta \tau \right) \right\rangle
=\left\vert b\right\rangle 
\end{align}%
This gives 
\begin{eqnarray}
x_{2}\left( t\right)  &=&x\left( t\right) -x\left( t\right) \delta \tau +%
\dot{x}\left( t\right) \left( \delta \tau \right) ^{2};\qquad \qquad t\in
\left( 0,\tau +\delta \tau \right]   \notag \\
x_{2}\left( t\right)  &=&b
\end{eqnarray}%
until we get to $\tau +\Delta \tau $, i.e. we continue the path in a sort of
analytic fashion. All other paths defined on the new time ineterval are then
obtianed by the increment functions of types 1 and 2.
\end{enumerate}

The general first order change $\left\vert \delta ^{\left( 1\right) }\Psi
\left( t\right) \right\rangle $ of a path $\left\vert \Psi \left( t\right)
\right\rangle \in \mathcal{P}\left( \infty ,\mathcal{H},\mathcal{H}\right) $
is then given by 
\begin{equation}
\left\vert \delta ^{\left( 1\right) }\Psi \left( t\right) \right\rangle
=\left\vert \Delta _{F}\Psi \left( t\right) \right\rangle +\left\vert \Delta
_{E}\Psi \left( t\right) \right\rangle +\left\vert \Delta _{\tau }\Psi
\left( t\right) \right\rangle 
\end{equation}%
where $\left\vert \Delta _{F}\Psi \left( t\right) \right\rangle $ is a first
order change that leaves the end points and time duration fixed, $\left\vert
\Delta _{E}\Psi \left( t\right) \right\rangle $ corresponds to a first order
change of the end points such that the end points are constrained to lie in
a manifold $\mathcal{M}\subseteq $ $\mathcal{H}$, while keeping the time
duration fixed and $\left\vert \Delta _{\tau }\Psi \left( t\right)
\right\rangle $ corresponds to a first order of the time duration, while
leaving the end points fixed. In order to obtain the most general
incremantal path $\left\vert \Delta _{F}\Psi \left( t\right) \right\rangle $
varies over the linear space $\mathcal{P}\left( \tau ,0,0\right) $ for a
fixed $\tau $, $\left\vert \Delta _{E}\Psi \left( t\right) \right\rangle $
varies over subspaces of paths described in type 2, and $\left\vert \Delta
_{\tau }\Psi \left( t\right) \right\rangle $ varies over the subspace of
paths described in type 4

\end{document}
